% =============================================================================
% Бүлэг 1: Онолын судалгаа
% =============================================================================

\chapter{Text2SQL-ийн онол ба судалгааны тойм}
\label{Chapter1}
\pagecolor{white}

\section{Text2SQL асуудал ба хөгжил}
\label{sec:t2s-problem}

Text2SQL даалгаврыг дараах байдлаар томьёолно: байгалийн хэлний асуулт $q$, өгөгдлийн сангийн схем $S$, нэмэлт мэдлэг
$K$ (баганын тайлбар, бизнес нэр томьёо, жишээ) өгөгдсөн үед $R(\hat{s}, D) = R(s^{\ast}, D)$ байх SQL $\hat{s}$-ийг
үүсгэх. Энд $s^{\ast}$ нь лавлах (gold) SQL, $R(s, D)$ нь $s$-ийг $D$ өгөгдлийн сан дээр гүйцэтгэсэн үр дүн. Нэг
асуултад олон зөв SQL байж болох тул текстийн тэнцүү байдлаар бус, гүйцэтгэлийн үр дүнгээр үнэлдэг.

Shi нар Text2SQL-ийн хөгжлийг дүрэмд суурилсан, Seq2Seq, урьдчилан сургасан хэлний загвар (PLM), LLM гэсэн дөрвөн
үе шатаар тодорхойлсон \cite{shi_2025_survey_llm_text2sql}. Энд PLM-ийг нейрон үед нэгтгэж гурван үеэр авч үзэв:
\begin{itemize}
  \item \textbf{Дүрэмд суурилсан систем.} Тодорхой домэйнд гараар бичсэн дүрэм. Шинэ схемд шилжүүлэх боломж бага.
  \item \textbf{Нейрон encoder--decoder.} Spider жишиг \cite{yu_2018_spider} гарснаар олон хүснэгттэй, урьд үзээгүй
        схемд ерөнхийлөх даалгавар тавигдсан. Схем холбох (schema linking), бүтэцтэй декодчилол хөгжсөн.
  \item \textbf{Том хэлний загвар.} Prompt-оор (DIN-SQL, DAIL-SQL), олон агенттай урсгалаар (MAC-SQL, CHASE-SQL), эсвэл
        тусгай fine-tune хийсэн загвараар (OmniSQL, Arctic-Text2SQL-R1) шийддэг
        \cite{pourreza_2023_din_sql, gao_2024_dail_sql, wang_2025_mac_sql, pourreza_2024_chase_sql, li_2025_omnisql,
        yao_2025_arctic_text2sql_r1}.
\end{itemize}

BIRD жишиг \cite{li_2023_bird} нь бодит, «бохир» өгөгдөл, том схем, гадаад мэдлэг (evidence) шаарддаг асуултаар
Spider-ээс илүү бодит нөхцөлийг загварчилсан. BIRD дээр evidence-ийг хасахад нарийвчлал огцом буурдаг нь домэйны мэдлэг
чухал болохыг харуулдаг. Энэ нь мэдлэгийн давхаргын (§\ref{sec:rag}) үндэслэл болно.

\section{Том хэлний загварт суурилсан архитектур}
\label{sec:llm-arch}

Орчин үеийн Text2SQL систем ихэвчлэн дараах бүрэлдэхүүнтэй дамжлага байдаг (Зураг~\ref{fig:generic-pipeline}):
асуултыг хэвийн болгох, мэдлэг хайх (RAG), SQL үүсгэх, шалгаж гүйцэтгэх, алдаа гарвал засах, олон нэр дэвшигч SQL-ээс сонгох.

\begin{figure}[htbp]
\centering
\begin{tikzpicture}[node distance=0.45cm and 0.5cm, every node/.style={font=\footnotesize},
  box/.style={draw, rounded corners, align=center, minimum height=0.9cm, text width=2.0cm, fill=gray!8}]
  \node[box] (q) {Асуулт\\(монгол)};
  \node[box, right=of q] (n) {Хэвийн\\болгох};
  \node[box, right=of n] (r) {Мэдлэг\\хайх (RAG)};
  \node[box, right=of r] (g) {SQL\\үүсгэх (LLM)};
  \node[box, below=0.9cm of g] (v) {Шалгах,\\гүйцэтгэх};
  \node[box, left=of v] (f) {Засах\\(алдаа)};
  \node[box, left=of f] (s) {Санал\\хураах};
  \node[box, left=of s] (a) {Хариулт\\эсвэл татгалзах};
  \draw[->] (q) -- (n); \draw[->] (n) -- (r); \draw[->] (r) -- (g); \draw[->] (g) -- (v);
  \draw[->] (v) -- (f); \draw[->] (f.north east) -- (g.south west); \draw[->] (f) -- (s); \draw[->] (s) -- (a);
\end{tikzpicture}
\caption{LLM-д суурилсан Text2SQL-ийн ерөнхий дамжлага}
\label{fig:generic-pipeline}
\end{figure}

\textbf{Prompt-д суурилсан аргууд.} DIN-SQL асуудлыг хэсэгчилж (schema linking, ангилал, үүсгэлт, өөрийгөө засах)
шийддэг \cite{pourreza_2023_din_sql}. DAIL-SQL нь асуултын утгыг placeholder-оор далдалж, төстэй асуулт ба SQL-ийн
бүтцээр жишээ сонгох нь few-shot-ийн чанарыг сайжруулдгийг харуулсан \cite{gao_2024_dail_sql}.

\textbf{Олон нэр дэвшигч ба сонголт.} CHASE-SQL нь олон генератор, сонгогч, засагчийг хослуулах нь ганц генератор
эсвэл энгийн санал хураалтаас илүү болохыг ablation-оор харуулсан \cite{pourreza_2024_chase_sql}. Хамгийн энгийн
хэлбэр нь хэд хэдэн нэр дэвшигч SQL-ийг гүйцэтгээд үр дүнгээр нь санал хураах (self-consistency) юм
\cite{wang_2023_self_consistency}. Санал нийлэлтийн хэмжээ нь итгэлийн хэмжүүр болж,
бага итгэлтэй үед хариулахаас татгалзахад ашиглагдана \cite{lee_2024_trustsql}.

\textbf{Тусгай загварууд.} OmniSQL нь 2.5 сая синтетик (асуулт, SQL) хосоор Qwen2.5-Coder-ийг сургаж, жижиг загварыг
том ерөнхий загваруудтай өрсөлдөхүйц болгосон \cite{li_2025_omnisql, hui_2024_qwen25_coder}. XiYan-SQL нь M-Schema
хэмээх хагас бүтэцтэй (semi-structured) схемийн дүрслэлийг санал болгосон \cite{gao_2024_xiyan_sql}.

\section{Хайлтаар баяжуулсан үүсгэлт (RAG)}
\label{sec:rag}

RAG нь загварын параметрт байхгүй мэдлэгийг хайлтаар олж prompt-д нэмдэг \cite{lewis_2020_rag}. Text2SQL-д дөрвөн
төрлийн мэдлэг хайдаг:
\begin{enumerate}
  \item \textbf{Схемийн тайлбар.} Хүснэгт, баганын тайлбар, синоним, enum кодын нэр. Багана нь англи нэртэй
        (\tsql{status\_code}) ч монгол хэрэглэгч «төлөв», «байдал» гэж асууна.
  \item \textbf{Утга.} Асуултад гарсан «Улаанбаатар», «мобайл апп» зэрэг үгийг санд хадгалсан утгатай
        (\tsql{'MOBILE'}, \tsql{'Улаанбаатар'}) тааруулах. CHESS нь LSH-ээр \cite{talaei_2024_chess}, CodeS нь BM25 ба
        LCS-ээр \cite{li_2024_codes} утга хайдаг.
  \item \textbf{Бизнес нэр томьёо.} «Идэвхтэй багц», «хугацаа хэтэрсэн зээл» гэх мэт байгууллагын тодорхойлолт. BIRD-ийн
        evidence-тэй ижил үүрэгтэй \cite{li_2023_bird}.
  \item \textbf{Жишээ.} Төстэй асуулт ба SQL (few-shot) \cite{gao_2024_dail_sql}.
\end{enumerate}

\textbf{Хайлтын арга.} Түлхүүр үгийн BM25 \cite{robertson_2009_bm25} нь яг тохирох нэр, кодыг сайн олдог бол
embedding-д суурилсан хайлт нь утгын төстэй байдлыг барьдаг. Хоёр арга бие биеэ нөхдөг тул тэдгээрийг жинлэн
хослуулах (hybrid) нь практикт түгээмэл. Олон хэлт BGE-M3 загвар 100 гаруй хэлийг дэмждэг ч уг өгүүлэлд монгол хэлийг
тусад нь үнэлээгүй \cite{chen_2024_bge_m3}.

\textbf{Schema linking хэзээ хэрэгтэй вэ.} Хүчирхэг загварын хувьд бүтэн схем контекстэд багтах үед schema linking
(хамаарах хүснэгтийг шүүх) нь хэрэгтэй хүснэгтийг алдах эрсдэлээс болж ашиггүй, бүр хортой байж болно. Гэвч уг
судалгаанд жижиг загварт (Llama~3.1-8B) schema linking EX-ийг өсгөсөн \cite{maamari_2024_death_schema_linking}.
Манай схем (28 view) контекстэд багтдаг тул энэ ажилд бүтэн схемийг анхдагч болгож, schema linking-ийн нөлөөг 7--9B
загвар дээр тусад нь (B10) хэмжинэ.

\section{Fine-tuning: SFT ба бататгах сургалт}
\label{sec:finetune}

\textbf{LoRA.} Загварын жинг бүхэлд нь шинэчлэхийн оронд сонгосон жингийн матрицуудад бага rank-тай засвар $\Delta W = BA$
($B \in \mathbb{R}^{d \times r}$, $A \in \mathbb{R}^{r \times k}$, $r \ll \min(d,k)$) нэмж сургадаг \cite{hu_2022_lora}.
Энэ нь 7--9B загварыг нэг GPU дээр сургах боломж олгоно. QLoRA нь суурь жинг 4 битээр шахдаг \cite{dettmers_2023_qlora}.

\textbf{Rejection sampling SFT.} Загвараар нэг асуултад олон хариу үүсгээд, гүйцэтгэлийн үр дүн нь gold-той таарсныг
нь л сургалтын өгөгдөлд үлдээнэ. Иймээс зөвхөн gold SQL-тэй (reasoning-гүй) асуултаас загварын өөрийн хэв маягтай, reasoning-тэй зөв жишээ олддог.

\textbf{GRPO.} Group Relative Policy Optimization нь нэг асуултад $G$ хариу (rollout) үүсгэж, reward-ийг бүлгийн
дунджаар хэвийн болгосон давуу талаар бодлогыг шинэчилнэ \cite{shao_2024_deepseekmath}:
\begin{equation}
A_i = \frac{r_i - \operatorname{mean}(r_1,\dots,r_G)}{\operatorname{std}(r_1,\dots,r_G)}, \qquad i = 1,\dots,G.
\label{eq:grpo}
\end{equation}
Утгын функцийн (value) тусдаа загвар шаардлагагүй тул санах ой хэмнэнэ. Arctic-Text2SQL-R1 нь зөвхөн гүйцэтгэлийн
үр дүнд суурилсан энгийн reward-аар (1: зөв, 0.1: ажилласан ч буруу, 0: алдаа) GRPO хийсэн. Түүний 32B хувилбар
BIRD-ийн нэг загварын самбарт 71.83\%-иар тэргүүлсэн бол энэ ажилд ашигласан 7B хувилбар BIRD test дээр 68.5\% EX
авсан \cite{yao_2025_arctic_text2sql_r1}. Олон хэлт GRPO-ийн судалгаа ч бага өгөгдлөөр мэдэгдэхүйц
ахиц гаргаж болохыг харуулсан \cite{kattamuri_2025_contrastive_grpo}.

\section{Agentic Text2SQL}
\label{sec:agentic}

Agent гэдэг нь загвар үйлдэл хийж (tool дуудаж), орчноос ажиглалт авч, дараагийн алхмаа шийдэх давталт юм. ReAct нь
бодол ба үйлдлийг ээлжлүүлдэг \cite{yao_2023_react}. Reflexion нь өмнөх оролдлогын алдаанаас дүгнэлт гаргаж дахин
оролддог \cite{shinn_2023_reflexion}. Text2SQL-д түгээмэл tool-ууд:
(1) судлах асуулга гүйцэтгэх (\tsql{SELECT DISTINCT channel ...});
(2) утга хайх;
(3) гүйцэтгэлийн алдааг буцааж засуулах;
(4) үр дүнг шүүмжлэгчээр (critic) шалгуулах.

Sqrl-9b нь \tsql{<sql>} (судлах) ба \tsql{<answer>} (эцсийн) хэмээх хоёр үйлдэлтэй agentic протоколоор сургагдсан
загвар \cite{feyn_2026_sqrl}. Харин LLM-ийн өөрийгөө засах чадвар хязгаарлагдмал: гадаад дохиогүй (гүйцэтгэлийн
алдаа, tool-ийн үр дүн) self-correction нь заримдаа зөв хариултыг эвддэг \cite{huang_2024_self_correct}. Энэ нь agent-ийн
хамгаалалтын дүрмийн (§\ref{sec:agent}) үндэслэл болно.

\section{Бага нөөцтэй хэл ба монгол хэлний сорилт}
\label{sec:lowres}

MultiSpider нь Spider-ийг 6 хэл рүү орчуулж (англитай нийт 7 хэл), англи бус хэл дээр нарийвчлал 6.1 нэгжээр
буурдгийг харуулсан \cite{dou_2023_multispider}. Spider 2.0-ийг 8 хэлд өргөтгөсөн MultiSpider 2.0 (өөр судлаачдын
ажил) дээр хамгийн сайн reasoning загварууд ч 4\% EX авсан бол MultiSpider 1.0 дээр 60\% байв
\cite{pham_2026_multispider2}. Бага нөөцтэй хэлэнд дараах стратегиуд туршигдсан:
\begin{itemize}
  \item \textbf{Translate-test.} Асуултыг англи руу орчуулах (гадны орчуулагч эсвэл загвар өөрөө) эсвэл загварыг
        англиар сэтгүүлэх нь англи бус оролттой харьцуулахад ихэвчлэн сайжруулдаг
        \cite{etxaniz_2024_self_translate, huang_2023_xlt}. Харин Text2SQL-д нэр, байршил зэрэг утга орчуулгаар гажих
        эрсдэлтэй \cite{dawar_2026_indicdb}.
  \item \textbf{Translate-train.} Англи сургалтын өгөгдлийг (BIRD, Spider) зорилтот хэл рүү орчуулж fine-tune хийх.
        BIRDTurk нь залгамал, SOV бүтэцтэй турк хэлэнд машин орчуулга ба хүний шалгалтын дамжлага боловсруулсан
        \cite{aktas_2026_birdturk}.
  \item \textbf{Олон хэлт жишээ.} Хэл хоорондын few-shot \cite{shi_2022_xricl}.
  \item \textbf{Hint ба галиг.} IndicDB нь латинаар бичсэн, англитай холилдсон (Hinglish) асуулт ба hint-ийн нөлөөг хэмжсэн
        \cite{dawar_2026_indicdb}.
\end{itemize}

MM-Eval нь кирилл монгол хэл дээр загваруудын утга зүй, сэтгэн бодох чадвар сул байгааг \cite{zhang_2024_mm_eval},
MonCulture-Eval нь монгол соёлын мэдлэгт бичиг хоорондын зөрүү байгааг \cite{minggad_2026_monculture_eval} харуулсан. Text2SQL-ийн
хувьд монгол хэлний гол бэрхшээлийг Хүснэгт~\ref{tab:mn-challenges}-д нэгтгэв.

\begin{table}[H]
\centering\small
\caption{Монгол хэлний Text2SQL-ийн онцлог бэрхшээл ба энэ ажлын хариу}
\label{tab:mn-challenges}
\begin{tabularx}{\textwidth}{>{\raggedright\arraybackslash}p{3.2cm}LL}
\toprule
Бэрхшээл & Жишээ & Энэ ажлын хариу \\
\midrule
Залгамал үгзүй & зээл, зээлийн, зээлээс, зээлүүдийн & Дагавар таслах дүрэм (F2), embedding хайлт \\
Кирилл ба латин галиг & «hvselt», «huselt», «хүсэлт» & Галигийн сул түлхүүр, A3/F2/F3 туршилт \\
Монгол утга санд & \tsql{city = 'Улаанбаатар'} & Утгын индекс (B2), литерал шалгах tool (G) \\
Код ба нэр & \tsql{'MOBILE'} = «мобайл апп» & Semantic model-ийн enum нэр (B1) \\
Бизнес нэр томьёо & «идэвхтэй багц», «PAR30» & Glossary (B3), хоёрдмол үгийг тодруулах \\
Загварын англи төвтэй байдал & 7B загварын монгол ойлголт сул & Орчуулга (A4), pivot (F5/F6), fine-tune (C) \\
\bottomrule
\end{tabularx}
\end{table}

\section{Үнэлгээ ба аюулгүй байдал}
\label{sec:eval-theory}

\textbf{Гүйцэтгэлийн нарийвчлал (EX).} Таамаг ба gold-ийн гүйцэтгэлийн үр дүн тэнцүү асуултын хувь:
\begin{equation}
\mathrm{EX} = \frac{1}{N}\sum_{j=1}^{N} \mathbf{1}\!\left[R(\hat{s}_j, D) = R(s^{\ast}_j, D)\right].
\label{eq:ex}
\end{equation}
EX нь санамсаргүйгээр таарсан (жишээ нь шүүлт буруу ч тоо нь давхацсан) таамгийг зөв гэж тооцох эрсдэлтэй.
\textbf{Test-suite нарийвчлал} нь өгөгдлийг өөрчилсөн хэд хэдэн хуулбар DB дээр зэрэг таарахыг шаарддаг
\cite{zhong_2020_test_suite}:
\begin{equation}
\mathrm{TS} = \frac{1}{N}\sum_{j=1}^{N} \prod_{i=0}^{m} \mathbf{1}\!\left[R(\hat{s}_j, D_i) = R(s^{\ast}_j, D_i)\right],
\qquad D_0 = D.
\label{eq:ts}
\end{equation}

\textbf{Найдвартай байдал.} TrustSQL нь хариулах боломжгүй асуултыг оруулж, зөв хариулт ба зөв татгалзалтад $+1$,
хариулах боломжтой асуултаас татгалзахад $0$, буруу хариултад $-c$ оноо өгдөг \cite{lee_2024_trustsql}:
\begin{equation}
\mathrm{Trust}_c = \frac{1}{N}\sum_{j=1}^{N}
\begin{cases}
+1, & \text{зөв хариулсан эсвэл зөв татгалзсан},\\
0, & \text{хариулах боломжтой асуултаас татгалзсан},\\
-c, & \text{буруу хариулсан эсвэл хариулах боломжгүйд хариулсан}.
\end{cases}
\label{eq:trust}
\end{equation}
TrustSQL өөрөө $c \in \{0, 10, N\}$-ийг ашигласан бол энэ ажилд санхүүгийн хэрэглээнд буруу тоо их хохиролтой гэж
үзэж $c \in \{1, 5, 10\}$-ийг сонгов.

\textbf{Хоёрдмол утга ба gold-ийн алдаа.} AMBROSIA, AmbiQT нь нэг асуулт олон зөв тайлбартай байж болохыг харуулсан
\cite{saparina_2024_ambrosia, bhaskar_2023_ambiqt}. Түгээмэл жишгүүдийн gold SQL-ийн нэлээд хэсэг алдаатай
байдаг тул eval set-ийг бие даан хянах шаардлагатай \cite{jin_2026_annotation_errors, wretblad_2024_bird_noise}.
Алдааг NL2SQL-BUGs-ийн ангиллаар шинжилдэг \cite{liu_2025_nl2sql_bugs}. LLM шүүгч (FLEX) нь EX-ийн дутагдлыг
нөхөх туслах хэмжүүр байж болно \cite{kim_2025_flex}.

\textbf{Статистик ач холбогдол.} Хоёр тохиргоог ижил асуулт дээр харьцуулахад McNemar-ын тест нь зөвхөн зөрсөн
хосыг ($n_{10}$: эхнийх зөв, хоёр дахь нь буруу; $n_{01}$: эсрэгээр) ашигладаг \cite{mcnemar_1947}. Жижиг түүвэрт яг
биномиал хувилбарыг хэрэглэнэ.

\textbf{Аюулгүй байдал.} LLM-ээр дамжсан SQL injection (prompt-to-SQL) нь асуултад заавар нуух замаар өгөгдөл
устгах, бусдын мэдээлэл уншихад хүргэж болно \cite{pedro_2025_p2sql}. Хамгаалалт нь давхаргатай байх ёстой: зөвхөн
унших эрхтэй DB хэрэглэгч, мөрийн түвшний хамгаалалт (RLS), SQL-ийн бүтцийн шалгалт, гүйцэтгэлийн хязгаар. Үйлдвэрлэлийн
системүүд ч ганц загвар биш, олон бүрэлдэхүүнтэй байдаг: Uber QueryGPT нь workspace, тусгай агентуудтай бөгөөд
бүрэлдэхүүн бүрийг тусад нь үнэлдэг \cite{uber_querygpt}. LinkedIn SQL Bot нь мэдлэгийн граф ашиглаж, асуулгыг
EXPLAIN-ээр шалгадаг \cite{linkedin_sqlbot}.

\section{Бүлгийн дүгнэлт}
\label{sec:ch1-summary}

Судалгааны тоймоос дараах дүгнэлтүүд гарав:
(1) Нээлттэй эхийн 7B орчим загварууд синтетик өгөгдөл ба RL-ийн тусламжтайгаар англи Text2SQL-д өндөр түвшинд
    хүрсэн. Гэвч англи бус хэл дээр мэдэгдэхүйц буурдаг.
(2) Домэйны мэдлэг (BIRD-ийн evidence) сургалтгүйгээр мэдэгдэхүйц нэмэр өгдөг. Харин schema linking-ийн ашиг загварын
    чадвар ба схемийн хэмжээнээс хамаардаг.
(3) Translate-train ба GRPO нь англи бус хэлэнд амжилттай туршигдсан ч монгол хэлэнд шалгагдаагүй.
(4) Agent ба self-correction нь гадаад дохиотой (гүйцэтгэл, tool) үед л найдвартай.
(5) Үнэлгээнд EX-ээс гадна test-suite, татгалзалт, хоёрдмол утга, статистик тест шаардлагатай.
Эдгээрт тулгуурлан дараагийн бүлэгт системийн зохиомжийг тодорхойлно.
